\section{Случайные блуждания и ветвящиеся процессы}

\subsection{Ветвящиеся процессы (Процесс Гальтона-Ватсона)}

\begin{example}[Модель размножения частиц]
Частицы делятся в каждый дискретный момент времени.
Пусть $X$ --- число потомков одной частицы, $P(X=k)=p_k$, $P_\xi(z) = \sum p_k z^k$.
Число частиц в поколении $t+1$:
$$ y_{t+1} = X_1 + X_2 + \dots + X_{y_t} $$
где $X_i$ независимы и одинаково распределены.
Производящая функция числа частиц:
$$ P_{y_{t+1}}(z) = P_{y_t}(P_\xi(z)) = \underbrace{P_\xi(P_\xi(\dots P_\xi(z)\dots))}_{t+1 \text{ раз}} $$
\end{example}

\begin{statement}[Математическое ожидание]
Найдем матожидание числа частиц в момент $t+1$:
$$ E[y_{t+1}] = P'_{y_{t+1}}(z)\Big|_{z=1} = P'_\xi(P_{y_t}(1)) \cdot P'_{y_t}(1)\Big|_{z=1} $$
Так как $P_{y_t}(1)=1$, получаем рекуррентное соотношение:
$$ E[y_{t+1}] = P'_\xi(1) \cdot E[y_t] = E[X] \cdot E[y_t] $$
Откуда $E[y_t] = (E[X])^t$ (при $y_0=1$).
\end{statement}

\begin{definition}[Вероятность вымирания]
Пусть $q_n = P(y_n = 0)$ --- вероятность смерти колонии (вымирания) к шагу $n$.
Так как событие $\{y_n=0\}$ влечет $\{y_{n+1}=0\}$, последовательность $q_n$ не убывает:
$$ 0 \le q_1 \le q_2 \le \dots \le 1 $$
Существует предел $\lim_{n \to \infty} q_n = q$.
\end{definition}

\begin{theorem}[Уравнение для вероятности вымирания]
Вероятность полного вымирания $q$ является наименьшим неотрицательным корнем уравнения:
$$ q = P_\xi(q) $$
\end{theorem}

\begin{proof}[Анализ решений]
Рассмотрим график $P(z)$ и прямую $y=z$.
\begin{itemize}
    \item \textbf{Вариант I:} $\mu = P'(1) \le 1$.
    В этом случае $P(z) > z$ для $z \in (0, 1)$. Единственный корень на $[0,1]$ --- это $z=1$.
    Следовательно, $q=1$ (вымирание достоверно).
    \item \textbf{Вариант II:} $\mu = P'(1) > 1$.
    Существует корень $\hat{q} \in (0, 1)$ такой, что $\hat{q} = P(\hat{q})$.
    В этом случае вероятность вымирания $q = \hat{q} < 1$.
\end{itemize}
\end{proof}

\subsection{Случайные блуждания (Возвратность)}

Рассмотрим простое случайное блуждание. Пусть $p$ --- вероятность шага вправо, $q=1-p$ --- влево.
Найдем вероятность $\tilde{p}$ того, что частица когда-либо вернется в начало координат (или достигнет уровня -1, начав с 0, что симметрично).
Из рекуррентного соотношения (по первому шагу):
$$ \tilde{p} = q + p \cdot \tilde{p}^2 $$
Решая квадратное уравнение $p\tilde{p}^2 - \tilde{p} + q = 0$:
$$ \tilde{p} = \frac{1 \pm \sqrt{1 - 4pq}}{2p} = \frac{1 \pm \sqrt{(p+q)^2 - 4pq}}{2p} = \frac{1 \pm |p-q|}{2p} $$
Корни: $\tilde{p}_1 = 1$ и $\tilde{p}_2 = \frac{q}{p}$.
Так как $\tilde{p}$ --- вероятность, то $\tilde{p} = \min(1, q/p)$.

\begin{statement}
Вероятность когда-либо вернуться в нуль (для симметричного или несимметричного случая):
$$ P(\text{возврат в } 0) = 1 - |p - q| $$
\end{statement}

\begin{example}[Производящая функция времени возвращения]
Для симметричного блуждания ($p=q=1/2$):
Пусть $Q(z)$ --- производящая функция момента первого возвращения.
$$ Q(z) = 1 - \sqrt{1 - z^2} $$
(В общем случае $Q(z) = 1 - \sqrt{1 - 4pqz^2}$).
\end{example}